\section{The audited problem and the result boundary}
\label{sec:scope}

The source of this continuation is the collective manuscript
\emph{Polylogarithms and their Arithmetic Bridges}, in
\src{Analysis/Polylogarithms/docs/manuscript} of ProveIt
\cite{Repo}. We fixed the repository revision before beginning the
mathematical audit. The revision is
\begin{center}\baseline.\end{center}
The reading artifact at that revision has twelve chapters and 375 pages.
The audit also inspected the associated research reports and the incoming
polylogarithm archives when determining which results had already been
obtained. The present deliverable is a separate article for integration;
it does not alter the repository remotely.

\subsection{What was already known}

Several attractive-looking targets were no longer open. In particular,
the $S_4$ evaluation, the stated Gaussian relation-space rank formulas,
the signed-measure threshold, principal slit-plane nonvanishing,
angular uniqueness, subcritical radial motion, fixed-total Gaussian
monotonicity, and the sharp critical Euler constant had already acquired
proofs. A previous continuation had also proved the first golden
weight-five identity. These facts determine the novelty boundary of
the present work: re-establishing one of them is useful only as a
clearly credited prerequisite or independent derivation.

The canonical manuscript still explicitly asks for a universal Euler
constant outside the critical triangle. It also gives a numerical
location at which the quartic radial coefficient vanishes along the
quadratic threshold, while leaving the higher local geometry open.
The two golden formulas with indices 20 and 24 are printed as numerical
identities without executable analytic certificates. The weight-seven
$S_6$ relation remains conjectural. These are the four interfaces
addressed here.

\begin{center}
\small
\begin{tabular}{@{}p{0.24\linewidth}p{0.32\linewidth}p{0.34\linewidth}@{}}
\toprule
Target & Baseline status & Contribution here\\
\midrule
Universal Euler error & Parameter-dependent supercritical bounds;
sharp critical constant & A universal $9/8$ bound for $N\ge2$;
the exact best constant for all $N\ge1$\\[5pt]
Golden index 20 and 24 & Numerical weight-five formulas & Exact
functional certificates and a proved finite descent\\[5pt]
Quartic transition & Existence and a diagnostic location & Exact local
uniqueness, negative sextic term, two turning points and scaling laws\\[5pt]
Cayley relations for $S_6$ & Abstract quotient and bounded search
obstructions & Explicit quotient projector and a nonzero full-ideal
normal form; the special-value identity remains open\\
\bottomrule
\end{tabular}
\end{center}

\subsection{Conventions}

For real $a,b>0$, put
\begin{equation}\label{eq:basic-family}
 H_m^{(b)}=\sum_{j=1}^{m}j^{-b},\qquad H_0^{(b)}=0,\qquad
 F_{a,b}(z)=\sum_{n=2}^{\infty}\frac{H_{n-1}^{(b)}}{n^a}z^n.
\end{equation}
The last series initially defines a holomorphic function on $|z|<1$.
For integer orders it is the depth-two polylogarithm
$\Li_{a,b}(z,1)$ in the nested-sum convention
$\sum_{n>m\ge1}z^n/(n^am^b)$. Values at $i$ mean the principal analytic
continuation, equivalently the radial Abel value. This convention matters
when $a+b\le1$, because the ordinary alternating coefficient series
need not converge. The positive double integral below defines the
continuation without applying a convergence theorem outside its domain.

We use $\beta(s)=\sum_{j\ge0}(-1)^j/(2j+1)^s$ and
$\eta(s)=\sum_{j\ge1}(-1)^{j-1}/j^s$ for $s>0$. Both series converge
on that domain. In the radial geometry, the conventional symbol
$\eta_{a,b}(\rho)$ denotes a normalized angle; its subscripts and
argument distinguish it from the Dirichlet eta function. To avoid a
collision between the Euler limit and a radial coefficient, the
quartic coefficient is always written with its arguments or in a
locally specified coordinate system.

The golden ratio and its reciprocal are
\[
 \phi=\frac{1+\sqrt5}{2},\qquad \rho=\phi^{-1},\qquad
 \ell=\log\rho=-\log\phi.
\]
The golden section uses this constant $\rho$; the radial section uses
$\rho$ as a variable radius and restates that convention locally.
All logarithms of positive real numbers are real logarithms. Real
parts of analytically continued polylogarithms at real arguments are
specified when intermediate functional identities leave $(0,1)$.

\subsection{What a certificate proves}

The Euler positivity verifier performs rational polynomial arithmetic
and proves positivity by Bernstein coefficients on a finite partition.
The radial verifier evaluates finitely many explicit elementary
expressions by outward rational interval arithmetic. The golden
certificates are exact tensor identities coupled to a written
differential and endpoint argument. The Cayley computation takes place
in a specified formal shuffle algebra. Its nonzero normal form proves
nonmembership in that algebraic relation ideal, and does not imply
nonvanishing of a numerical period.

These are ordinary mathematical proofs with finite computational
premises. The executable code and exact records expose those premises
for review. Numerical quadrature, plots, and stationary-point solvers
serve only as diagnostics. They are neither the source of an exact
sign claim nor a substitute for the analytic arguments.
